\section{Discussion}
\label{sec:discussion}

This section interprets the results of Section~\ref{sec:results} in the wider HTR and interpretable-ViT context. It is organized around the four findings the experiment matrix produced, each read as a claim about the CTC-HTR setting rather than a claim about the specific runs that produced it.

\subsection{Register Tokens Change Attention Structure Without Changing Recognition}
\label{sec:disc_registers}

The register-token intervention proposed by Darcet~\etal~\cite{darcet2024registers} was designed and validated on classification ViTs, where the accuracy-neutrality claim and the attention-cleanliness claim are two sides of the same experiment: cleaner patch attention improves downstream detection and segmentation but leaves ImageNet top-1 alone. Section~\ref{sec:res_claim} of this report shows that the two claims separate in the CTC HTR setting. The accuracy-neutrality claim carries over cleanly: within a seed-noise floor of $\pm 0.08$~pp on IAM and $\pm 0.09$~pp on READ2016, the recognition error rate is flat across $R \in \{0, 2, 4, 8, 16\}$. The attention-cleanliness claim carries over partially: attention entropy decreases at $R \in \{4, 8\}$, the Spearman correlation between character position and attention peak improves, and the number of attention heads aligned to left-to-right reading order grows monotonically with $R$; but the token-norm asymmetry that would signal registers absorbing sinks in a classification ViT is inverted here --- patch tokens remain the higher-norm carriers, and a residual secondary lobe survives on those patch columns across every character in the per-character attention maps.

The mechanistic explanation given in Section~\ref{sec:res_claim} is that a CTC self-attention encoder has no CLS-style dedicated query. The tokens that produce the CTC-decoded output are the patch tokens themselves, and if a patch token has developed high norm, it is both a query and a key of the sink pattern. Registers cannot dissolve this self-consistency from one side alone. This is a structural difference between the two settings and it is a testable prediction: adding a CLS-style aggregate query on top of the ViT-RGTS~v2 encoder should let registers take over the global-information role more completely, closing the residual gap. The Supplementary document sketches this as the first extension.

The practical takeaway for HTR practitioners is straightforward. If the objective is minimum CER on IAM-scale data, register tokens are neither a benefit nor a harm --- the extra $R \cdot D$ parameters are essentially free but do not move the recognition needle. If the objective is interpretable attention for a downstream task (writer analysis, style analysis, character-indexed features for retrieval), $R = 4$ or $R = 8$ is the sweet spot on the entropy and ordering metrics reported in Fig.~\ref{fig:register_quant}, and the residual sink pattern should be suppressed at inference time via a CTC-column prior on the raw attention, as documented in Section~\ref{sec:attn_eval}.

\subsection{The CNN Stem Is Not a Refinement, It Is a Structural Requirement}
\label{sec:disc_cnnstem}

Two independent code paths --- ViT-RGTS~v1 with row-major patch flattening, and ViT-RGTS~v2 with \texttt{use\_cnn\_stem=False} --- both collapse to ${\sim}72$--$74\%$~CER (Section~\ref{sec:res_abl_arch}, Table~\ref{tab:ablations}). Every other ablation in the experiment matrix moves the CER by at most a few percentage points; these two move it by $67$~pp. The mechanism is not architectural finesse or hyperparameter sensitivity; it is that CTC's monotonic-alignment assumption requires the token sequence entering the CTC head to be strictly left-to-right in reading order, and a row-major flattening of a 2D patch grid interleaves vertical position into the sequence order, so the CTC forward--backward algorithm has no valid alignment between the ground-truth character sequence and the emitted token order.

The consequence for CTC ViT HTR architects is that the token-formation step cannot be treated as a preprocessing detail. Whichever mechanism collapses the 2D patch grid to a 1D reading-order sequence --- a CNN stem with height-collapsing pooling, an adaptive max-pool as in HTR-VT~\cite{li2024htrvt}, or an explicit column-attention pool --- is the load-bearing element. Register tokens, model width, and even model depth are secondary: none of these ablations move the CER by more than ${\sim}1$~pp, but a wrong token order alone destroys $67$~pp. Two convergent failures in this report isolate the causal pathway; a third exists in the historical HTR-VT development, which independently arrived at the same design decision~\cite{li2024htrvt}.

\subsection{SWA and Fine-Tuning Strategy Are Orthogonal, and Both Matter}
\label{sec:disc_swa_ft}

Two training-side levers move CER independently. Stochastic Weight Averaging~\cite{izmailov2018swa} adds a $0.26$--$0.39$~pp improvement to any run whose loss trajectory has genuinely settled into a flat basin by the SWA start epoch; the effect is consistent across CNN--RNN, ViT-RGTS~v2, and TrOCR (Section~\ref{sec:res_finetune}). SWA hurts, sometimes catastrophically, when the underlying run has not converged --- averaging over a wandering trajectory does not produce a better weight vector, and the SWA of the naive-fine-tune ViT-B/16 run is worse ($+16.96$~pp) than its non-SWA counterpart. The signal for whether SWA will help is available from the training log itself: if the CTC loss has plateaued and the CER curve has stopped improving materially in the fifteen epochs before the SWA start, SWA is safe to enable; otherwise it is not.

The fine-tuning strategy choice for pretrained models is a second, independent lever. The $73.14\% \to 15.09\%$~CER improvement from naive uniform LR to two-group LR on ViT-B/16 (Section~\ref{sec:res_finetune}, Table~\ref{tab:finetune}) is entirely attributable to the LR grouping decision. The strategy that wins depends on the pretrained-representation--target-domain gap: two-group LR (uniform low backbone rate, high head rate) wins for large domain gaps (ImageNet$\to$handwriting), LLRD wins for small domain gaps (TrOCR encoder pretrained on handwriting), and gradual unfreezing fails in both regimes at the epoch budgets tested. The pattern generalizes across pretrained-transfer applications, but the specific LR ratio and decay factor should be re-tuned per model family.

\subsection{Synthetic Handwriting Is Not a Free Substitute for IAM}
\label{sec:disc_synth}

The synthetic-pretraining lane (Section~\ref{sec:res_synth}) is the most compressed lesson in the matrix. Models reach $0.5$--$0.8\%$~CER on the synthetic validation set --- effectively perfect --- but $30$--$46\%$~CER on real IAM test. That gap is the render-to-real distribution shift: no ink bleed, no writer-specific idiosyncratic letterforms, no paper texture, no scanning artifacts appear in the rendered corpus, and models trained exclusively on it fit the rendering-specific regularities rather than the underlying character-shape structure that transfers. The Nikolaidou~\etal~\cite{nikolaidou2023wordstylist} pipeline does exactly what it was designed to do --- render clean, controlled handwriting-style text at scale --- but its outputs are not a drop-in substitute for IAM training data for a from-scratch encoder.

The observation that ViT-RGTS~v2 transfers slightly better than CNN--RNN from synthetic to real ($30.31\%$ vs.\ $35.45\%$ CER at $R = 0$) suggests --- but does not prove --- that self-attention is architecturally less prone to overfitting the pixel-level rendering signature than a CNN filter bank. This is consistent with the broader literature on ViT domain generalization but is not tightly controlled in this report, and the walltime truncation on the synthetic runs prevents a converged conclusion. A curriculum that pretrains on synthetic and fine-tunes on IAM is the natural follow-up experiment, and would likely close a substantial fraction of the gap.

\subsection{Limitations}
\label{sec:disc_limits}

Three limitations bound the claims of this report.

\paragraph{Attention quality is measured, not evaluated end-to-end.} The interpretability half of the Darcet claim is supported here by three attention-side metrics (entropy, spatial ordering, aligned heads) and one qualitative reproduction (BM per-character maps). It is not evaluated by measuring downstream writer identification accuracy on top of the attention maps, because a full writer-ID experiment is out of scope for the 10~ECTS project budget. The Supplementary document sketches how a VLAC-style pipeline~\cite{souibgui2022vlac} would plug into the ViT-RGTS~v2 attention output; running that experiment is the single most valuable follow-up.

\paragraph{Scale.} Every claim is measured at IAM scale ($6{,}482$ training lines) and READ2016 scale ($8{,}349$ training lines). It is possible that at $10 \times$ more training data the CER-vs-$R$ curve becomes non-flat, or that at the multi-million-line scale used by TrOCR the register mechanism starts to show accuracy gains. The synthetic-pretraining lane hinted at this direction but was walltime-limited.

\paragraph{One CTC decoder.} All CER numbers use greedy best-path CTC decoding. Beam-search decoding with a language model would lower every reported CER but is unlikely to change the register-count \emph{ordering}. The choice to skip it here is scope, not principle: comparing a language-model-augmented ViT-RGTS~v2 against a language-model-augmented CNN--RNN would be a stronger absolute claim but would not affect the register-vs-CER question this report was set up to answer.
